\documentclass[10pt,a4paper]{amsart}
\usepackage[margin=19mm]{geometry}
\usepackage{amsmath,amssymb,mathtools}
\usepackage[scr=rsfs,cal=euler]{mathalfa}
\usepackage{xfrac}
\usepackage[unicode,hidelinks,bookmarks=false]{hyperref}
\newcommand{\ie}{i.e.\ }
\newcommand{\cf}{cf.\ }
\newcommand{\resp}{respectively\ }
\newcommand{\Subsection}{\subsection}
\providecommand{\nobreakdash}{\nobreak\hbox{-}}
\setlength{\emergencystretch}{2em}
\multlinegap0pt
\usepackage{bm}
\usepackage{bbm}
\usepackage{dsfont}
\usepackage{xspace}
\usepackage{cancel}
\usepackage{mathtools}
\usepackage{cleveref}
\usepackage{amsthm}
\makeatletter
\@ifundefined{theorem}{\newtheorem{theorem}{Theorem}}{}
\@ifundefined{prop}{\newtheorem{prop}[theorem]{Proposition}}{}
\makeatother
\DeclareMathOperator{\rank}{rank}
\newcommand{\R}{\mathbb{R}}
\newcommand{\paramspace}[1]{\Theta_{#1}}
\title[The Symmetries of Three-Layer ReLU Networks]{The Symmetries of Three-Layer ReLU Networks}
\author{Johanna Marie Gegenfurtner, Moritz Grillo, Guido Mont\'ufar}
\date{Source: arXiv:2605.18319v1 (CC BY 4.0)}
\begin{document}
\maketitle

\noindent\emph{Source and licence.} The Symmetries of Three-Layer ReLU Networks. Version arXiv:2605.18319v1 (CC BY 4.0), distributed under the Creative Commons Attribution 4.0 International licence, as recorded in the arXiv licence metadata for this exact version and shown on the abstract page https://arxiv.org/abs/2605.18319v1. Full rights notice: licenses/arxiv-2605.18319-CC-BY-4.0.txt.\par\smallskip

\noindent\textit{Editorial scope.} Counted unit: the Prop whose source label is \texttt{prop:dim\_special\_case} in the article (source0.tex lines 2061--2067, source label \texttt{prop:dim\_special\_case}), delivered together with the article's own complete original proof of it (source0.tex lines 2069--2155, one \texttt{proof} environment of 87 lines). No mathematics is written, repaired or re-derived: statement and proof are the article's own text line for line. Required context retained uncounted: prop:empty\_or\_dim (lines 1956--1959). Every definition, display and label of those lines is retained unchanged. Omitted from this package and not redistributed: 1--1955, 1960--2060, 2068--2068, 2156--3474. Nothing inside a retained excerpt is omitted or altered: every retained body is delivered exactly as the article's lines are written, so no omission receipt of any kind accompanies this package. Citations: the retained excerpts carry no citation command at all, so no bibliography entry of the article is transcribed and no reference is redistributed. Reference adaptation: none was needed and none was made; every reference inside the delivered unit resolves inside it, so no unresolved reference or citation is produced. Printed slips kept literal: no printed word, symbol, hypothesis or number of the counted statement or of its proof was altered. Layout and class: the article's own class is \texttt{article} with packages including myarxiv, inputenc, fontenc, hyperref, url, booktabs, amsfonts, nicefrac, microtype, xcolor, graphicx, amsmath, amssymb, amsthm; the wrapper is the lane's pinned \texttt{amsart} editorial wrapper at 10pt on A4, which restates the theorem-like environments the retained text needs and adds the article's own macro definitions that the retained text needs (\texttt{\textbackslash R}, \texttt{\textbackslash paramspace}, \texttt{\textbackslash rank}), with extra wrapper packages (bm, bbm, dsfont, xspace, cancel, mathtools, cleveref). The delivered numbering is the unit's own. Assets: no retained span contains a figure environment, a graphics-inclusion command, a PDF-page inclusion, a table or a TikZ picture; no image file is copied into this package, the package declares no asset dependency and no image file is redistributed with it. Licence: version arXiv:2605.18319v1 of this article is distributed under the Creative Commons Attribution 4.0 International licence, as recorded in the arXiv licence metadata for this exact version and shown on the abstract page https://arxiv.org/abs/2605.18319v1. A full rights notice is delivered at licenses/arxiv-2605.18319-CC-BY-4.0.txt. Characters: every retained excerpt is pure ASCII, so the delivered engine, XeLaTeX with its default Latin Modern text font, reports no missing character.
\begin{prop}
\label{prop:empty_or_dim}
Let $\theta \in \paramspace{(d,n,m)}$ be generic and minimal and let $k$ denote the number of nonlinear neurons. If with $n-k \neq m+1$, then, for every $S \subseteq [k]$, either $\dim(K_S) = (n-k)^2$ or $K_S=\emptyset$. 
\end{prop}

\begin{prop}
\label{prop:dim_special_case}
Let $\theta \in \paramspace{(d,n,m)}$ be generic and minimal and let $k$ denote the number of nonlinear neurons. If $n-k = m+1$, then, for every $S \subseteq [k]$, we have
$
\dim(K_S)= m^2 +m +d. 
$
\end{prop}

\begin{proof}
Let $M_S$ be as in the proof of \Cref{prop:empty_or_dim}. Since $\theta$ is minimal, we have $d \geq m+1$. By genericity, $\rank(A_S) = m$. 

We first compute the dimension of
$
\overline{M_S}
=
\{(U,V)\in \R^{m\times (m+1)}\times \R^{(m+1)\times d}\mid UV=A_S\}.
$
Fix $U \in \R^{m \times (m+1)}$ with $\rank(U)=m$. Then $\dim\ker(U)=1$, so for each column $a_j$ of $A_S$, the solution set of $Ux=a_j$ is an affine line. Hence
$
\{V \in \R^{(m+1) \times d} \mid UV =A_S\}
$
is an affine space of dimension $d$. Therefore
$
\dim(\overline{M_S}) = m(m+1)+d= m^2 + m +d.
$

Now let
$
Z= \{ U \in \R^{m \times (m+1)} \mid \exists v \in \R^{m+1}_{>0}:  \ker(U) = \operatorname{span}(v)\}.
$
This is an open subset of $\R^{m\times (m+1)}$, hence
$
\dim(Z)=m(m+1)=m^2+m.
$

For $U \in Z$, define
$
M_S(U) \coloneqq \{V \in \R^{(m+1) \times d} \mid UV =A_S ,\ V_{i,j} \geq 0\}.
$
We claim that $\dim(M_S(U))=d$. Indeed, since $\rank(U)=m$, the affine space
$
\{V \in \R^{(m+1) \times d} \mid UV =A_S\}
$
has dimension $d$. Choose $c>0$ sufficiently large so that
$
Y_c = \{V \in \R^{(m+1) \times d} \mid UV =A_S ,\ |V_{i,j}| \leq c\}
$
has nonempty interior relative to this affine space; in particular, $\dim(Y_c)=d$.

Let $v\in \R^{m+1}_{> 0}$ satisfy $\ker(U) = \operatorname{span}(v)$ and set
$
\delta = \min_i v_i >0.
$
Define
$
\varphi_{c,v}\colon
\begin{bmatrix}
V_1 & \cdots & V_d
\end{bmatrix}
\mapsto
\begin{bmatrix}
V_1 + \frac{c}{\delta}v & \cdots & V_d +\frac{c}{\delta}v
\end{bmatrix}.
$
Since $Uv=0$, the map $\varphi_{c,v}$ preserves the equation $UV=A_S$. Moreover, for every entry we have
$
(V_j)_i + \frac{c}{\delta}v_i \ge -c + \frac{c}{\delta}\delta =0,
$
so $\varphi_{c,v}(Y_c) \subseteq M_S(U)$. Since $\varphi_{c,v}$ is an affine isomorphism, it follows that
$
\dim(M_S(U)) \ge d.
$
On the other hand,
$
M_S(U)\subseteq \{V\in \R^{(m+1)\times d}\mid UV=A_S\},
$
and the latter has dimension $d$, hence $\dim(M_S(U))=d$.

Therefore
$
M_S\cap (Z\times \R^{(m+1)\times d})
$
projects onto $Z$ with all fibers of dimension $d$. By the fiber dimension theorem for semi-algebraic sets,
$
\dim(M_S)\ge m^2+m+d.
$
Since $M_S\subseteq \overline{M_S}$ and $\dim(\overline{M_S})=m^2+m+d$, we conclude that
$
\dim(M_S)=m^2+m+d.
$

Finally, passing from $M_S$ to $K_S$ is analogous to the proof of \Cref{prop:empty_or_dim}: the translation variables contribute $m+1$ degrees of freedom, while the row-normalization constraints remove exactly $m+1$ degrees of freedom, and the output bias is then uniquely determined. Hence 
$\dim(K_S)=\dim(M_S)=m^2+m+d$. 
%
\end{proof}

\end{document}
